\documentclass[sigconf,screen]{acmart}
\setcopyright{acmlicensed}
\copyrightyear{2025}
\acmYear{2025}
% TODO before submission: replace with actual conference name, date, and city
\acmConference[CONFERENCE'25]{ACM Conference}{Month, 2025}{City, Country}
\acmBooktitle{Proceedings of CONFERENCE'25}
\acmDOI{10.1145/XXXXXXX.XXXXXXX}
\acmISBN{978-1-4503-XXXX-X/25/XX}

\usepackage{booktabs}
\usepackage{multirow}
\usepackage{amsmath}
\usepackage{graphicx}
\usepackage{subcaption}
\usepackage{siunitx}
\usepackage{tabularx}
\usepackage{enumitem}
\usepackage{xspace}
\usepackage{url}
\usepackage{flushend}
\usepackage{float}

\newcommand{\method}{GenomeRL\xspace}
\newcommand{\dnabert}{DNABERT-2\xspace}
\newcommand{\comp}{\mathrm{Comp}}
\newcommand{\tokens}{\mathrm{Tok}}

\title{\method: Reinforcement Learning for Adaptive DNA Tokenization in Genomic Sequence Classification}

\author{Baran Khazaee}
\affiliation{%
  \institution{University of California Davis}
  \city{Davis}
  \country{USA}
}
\email{fkhazaee@ucdavis.edu}

\author{Michael Cho}
\affiliation{%
  \institution{University of California Davis}
  \city{Davis}
  \country{USA}
}
\email{mhccho@ucdavis.edu}

\author{Xinyuan Jiang}
\affiliation{%
  \institution{University of California Davis}
  \city{Davis}
  \country{USA}
}
\email{xypjiang@ucdavis.edu}

\settopmatter{printacmref=false}
\renewcommand\footnotetextcopyrightpermission[1]{}

\begin{document}

\begin{abstract}
Genomic foundation models rely on fixed tokenization schemes that impose uniform segmentation regardless of local information density. We study adaptive DNA tokenization as a reinforcement-learning problem in which a boundary policy learns sequence segmentations that trade off predictive performance against compression, and critically evaluate it against simpler non-RL baselines. Experiments span a prototype end-to-end promoter model, a controlled ablation, and proposal-faithful frozen and fine-tuned \dnabert~environments, with comparisons to fixed-stride mean-pooling and Gumbel-Softmax differentiable segmentation at matched compression budgets. A central finding is that a simple fixed-stride pooling baseline is competitive with or better than RL-learned boundaries on most tasks (F1 0.920 vs.\ 0.886 on promoter; near-baseline 0.925 vs.\ 0.929 on splice-site), suggesting that compression itself — rather than learned boundary placement — is the primary driver of performance in co-adapted settings. The RL approach's distinctive value lies in learning biologically meaningful boundary structure: boundaries are specifically enriched around TATA-box elements but not other regulatory motifs, a pattern absent in fixed-stride pooling. Frozen-backbone RL consistently over-compresses and fails, while fine-tuned RL achieves favorable compression–performance tradeoffs on TF-binding tasks. The results establish compression-aware tokenization as a viable research direction while identifying boundary-sensitive regularization, constrained optimization, and segment-aware positional encoding as key open challenges.
\end{abstract}

\keywords{DNA tokenization, reinforcement learning, genomic language models, adaptive segmentation, DNABERT-2, sequence compression}

\maketitle
\pagestyle{plain}

\section{Introduction}
DNA language models have become a central tool for genomic prediction, but nearly all current pipelines depend on fixed tokenization. In practice, this means that every region of a sequence is segmented at the same granularity, regardless of whether that region is highly informative, repetitive, or biologically structured. This uniform treatment is convenient, but it may be suboptimal for tasks in which different subsequences require different levels of resolution.

This paper investigates whether tokenization itself can be learned adaptively. We frame DNA segmentation as a sequential decision problem in which an RL policy places token boundaries along a nucleotide sequence. The policy receives feedback from a downstream classifier together with a compression-oriented reward, allowing it to discover segmentations that balance predictive quality against representation efficiency.

We study this question on promoter detection and then broaden the evaluation to additional GUE tasks. The work makes four contributions:
\begin{enumerate}[leftmargin=*,nosep]
\item We implement a prototype \method system that jointly studies learned segmentation and downstream prediction on promoter detection.
\item We perform controlled tokenizer ablations that isolate the effect of segmentation while keeping the classifier fixed.
\item We build proposal-faithful frozen and fine-tuned \dnabert environments to test whether RL tokenization remains useful when paired with stronger pretrained genomic encoders.
\item We analyze interpretability and cross-task transfer, including five transcription-factor binding tasks and one splice-site task from GUE.
\end{enumerate}

The results reveal a more complex picture than originally anticipated. A fixed-stride mean-pooling baseline at comparable compression outperforms fine-tuned RL on promoter detection (F1 0.920 vs.\ 0.886) and nearly matches the full \dnabert baseline on splice-site prediction (0.925 vs.\ 0.929), while RL collapses to 0.241 on the same task. This suggests that compression-based regularization — not learned boundary placement — drives most of the performance benefit in co-adapted settings. The RL approach's distinctive contribution is biological interpretability: learned boundaries are specifically enriched around TATA-box elements, a pattern absent from stride pooling. Frozen-backbone RL systematically over-compresses and fails. These findings sharpen the scope of adaptive tokenization as a research direction and identify constrained optimization, segment-aware positional encoding, and task-aware reward design as key open challenges.

\section{Related Work}

\subsection{Genomic Foundation Models}
Transformer language models \cite{vaswani2017attention,devlin2019bert,liu2019roberta} have been adapted to DNA sequence modeling across a range of architectures and scales. DNABERT \cite{ji2021dnabert} pioneered k-mer-based pretraining on the human genome; DNABERT-2 \cite{zhou2024dnabert2} replaced k-mers with a byte-pair-encoded vocabulary and introduced the GUE benchmark for multi-task evaluation. Long-context architectures have extended genomic modeling to much longer sequences: HyenaDNA \cite{nguyen2023hyenadna} uses sub-quadratic convolution operators to model single-nucleotide sequences up to one million bases, while Enformer \cite{avsec2021effective} integrates long-range interactions for gene expression prediction. The Nucleotide Transformer \cite{dallatorre2023nucleotide} scales BERT-style pretraining to multi-species genomes, and Caduceus \cite{schiff2024caduceus} introduces a bidirectional state-space architecture for DNA. Evo \cite{nguyen2024sequence} extends sequence modeling to genome-scale design across molecular to genomic contexts. Despite this diversity, \emph{all} of these models fix their tokenization scheme at pretraining time and leave it unchanged at the downstream task — the core limitation our work addresses.

\subsection{Tokenization in Language and Biology}
Subword tokenization methods such as Byte-Pair Encoding \cite{sennrich2016neural} and SentencePiece \cite{kudo2018sentencepiece} learn data-driven vocabularies from text corpora and have become standard in NLP. In genomics, the analogous choices are single nucleotides, fixed k-mers, or pretrained BPE vocabularies. All three are static: once chosen, the segmentation granularity does not adapt to downstream task structure. Our work asks whether a \emph{task-adaptive} tokenizer can improve the compression-performance tradeoff compared to these fixed alternatives.

\subsection{Regulatory Sequence Analysis}
Deep learning models for regulatory genomics have achieved strong performance on promoter detection, splice-site prediction, and non-coding variant effect prediction \cite{zhou2015predicting,kelley2018sequential,jaganathan2019predicting}. SpliceAI \cite{jaganathan2019predicting} demonstrated that deep networks can learn splice acceptor and donor signals from raw sequence, establishing a strong baseline for boundary-sensitive prediction that our transfer experiments revisit. Motif-based interpretability — comparing model behavior on known transcription factor binding sites from databases such as JASPAR \cite{fornes2020jaspar} — is a standard diagnostic for genomic models \cite{avsec2021effective}, and we adopt this approach to analyze the learned tokenizer.

\subsection{Reinforcement Learning for Discrete Decisions}
Policy-gradient methods optimize non-differentiable discrete decisions by estimating the gradient of expected reward \cite{williams1992simple,sutton2018reinforcement}. REINFORCE \cite{williams1992simple} provides the foundational estimator we build on, and Proximal Policy Optimization \cite{schulman2017proximal} offers a clipped surrogate that stabilizes training — a direction we identify as future work. RL has been applied to discrete text generation in NLP: SeqGAN \cite{yu2017seqgan} trains a sequence generator with a discriminator reward, and self-critical sequence training (SCST) \cite{rennie2017self} optimizes non-differentiable captioning metrics. Both use REINFORCE-style gradients over discrete token sequences, analogous to our boundary decisions. Closer to our problem, RL has been used for word segmentation and morpheme boundary detection in NLP \cite{yang2018toward}: these methods learn to split or merge character sequences to improve downstream task performance, sharing the core idea of task-driven tokenization but operating over natural language rather than DNA and without an explicit compression objective. Adaptive computation approaches including Adaptive Computation Time \cite{graves2016adaptive} and depth-adaptive transformers \cite{elbayad2020depth} vary computation per input at the depth dimension rather than the tokenization granularity. Token merging \cite{bolya2023token} and dynamic token pooling \cite{nawrot2022dynamic} reduce sequence length in vision and NLP transformers by merging similar tokens during attention — a post-encoding analogue of our pre-encoding boundary decision. Differentiable segmentation via Gumbel-Softmax relaxations \cite{jang2017categorical} offers a gradient-based alternative to discrete boundary sampling; unlike our REINFORCE estimator, such approaches backpropagate directly through the segmentation but require a relaxation temperature schedule and do not naturally encode a compression budget. Length-adaptive transformers and token pruning methods \cite{kim2022learned} drop uninformative tokens at intermediate layers to reduce FLOPs; these operate on internal representations rather than the input tokenization stage and do not produce interpretable segment boundaries. Recent work in vision-language models (VLMs) has explored budget-aware routing that explicitly separates low-frequency ``anchor'' tokens from high-frequency ``detail'' tokens under compression budgets \cite{chen2025parcel}, showing that frequency-role specialization can shift the performance–efficiency frontier on fine-grained tasks. Our mean-pooling approach lacks an analogous mechanism for preserving high-frequency boundary signals, which helps explain the splice-site failure; a two-branch segment-anchor design is a promising future direction. Our work, to our knowledge, is the first to apply policy-gradient RL specifically to the tokenization step of a genomic classification pipeline.

\section{Method}
\subsection{Adaptive DNA tokenization as RL}
Let a DNA sequence be $x=(x_1,\ldots,x_L)$ with length $L$. At each position $t$, a policy $\pi_\theta$ predicts a binary boundary decision
\begin{equation}
    b_t \in \{0,1\}, \qquad b_t \sim \pi_\theta(\cdot \mid s_t),
\end{equation}
where $s_t$ denotes the local state available to the boundary agent at position $t$.

The boundary vector $b=(b_1,\ldots,b_L)$ induces a segmentation
\begin{equation}
    z = S(x,b) = (z_1,\ldots,z_{T(b)}),
\end{equation}
where each $z_i$ is a contiguous subsequence of $x$ and $T(b)$ is the resulting number of tokens. We define
\begin{equation}
    T(b) = 1 + \sum_{t=2}^{L} b_t,
\end{equation}
and the normalized compression ratio
\begin{equation}
    C(b) = \frac{T(b)}{L}.
\end{equation}
\textbf{Compression ratio denominator.} Throughout this paper, $C(b) = T(b)/L$ is always normalized by $L$, the number of \emph{\dnabert BPE tokens} in the original (ungrouped) sequence — not by raw nucleotide count. For example, a 300\,bp promoter produces $L \approx 63$ BPE tokens; compression $C = 0.19$ means $\approx 12$ segments. All tables use this same BPE-token denominator consistently. The BPE token count varies slightly across sequences due to variable-length BPE merges, so per-sequence compression ratios have small variance; reported values are means over the test set.

The segmented sequence is passed to a downstream classifier $f_\phi$, which produces a predictive distribution
\begin{equation}
    p_\phi(y \mid z) = p_\phi(y \mid S(x,b)).
\end{equation}

We optimize a reward that trades off task loss against compression:
\begin{equation}
    R(x,y,b) = -\mathcal{L}_{\mathrm{cls}}\!\left(y, p_\phi(y \mid S(x,b))\right) - \beta C(b),
\end{equation}
where $\beta > 0$ controls the strength of the compression penalty.

The tokenization policy objective is
\begin{equation}
    J(\theta) = \mathbb{E}_{b \sim \pi_\theta(\cdot \mid x)}[R(x,y,b)].
\end{equation}
We optimize this objective with a policy-gradient estimator of REINFORCE form \cite{williams1992simple,sutton2018reinforcement}:
\begin{equation}
    \nabla_\theta J(\theta) \approx
    \left(R - \hat{R}\right)
    \sum_{t=1}^{L}\nabla_\theta \log \pi_\theta(b_t \mid s_t),
\end{equation}
where $\hat{R}$ is a baseline used to reduce variance.

In the fine-tuned setting, the classifier parameters $\phi$ are also updated by minimizing the supervised loss:
\begin{equation}
    \min_{\phi}\;\mathbb{E}_{b \sim \pi_\theta}[\,\mathcal{L}_{\mathrm{cls}}(y, p_\phi(y \mid S(x,b)))\,].
\end{equation}

\subsection{Overall workflow}
Figure~\ref{fig:method} summarizes the overall workflow used in this study. We begin with a DNA sequence, apply an RL boundary policy to produce grouped tokens, and score the grouped representation with a downstream classifier. The same mechanism is then evaluated across four stages: a prototype promoter model, a controlled tokenizer ablation, proposal-faithful frozen and fine-tuned \dnabert environments, and finally interpretability plus transfer experiments on GUE tasks. This staged design lets us move from proof of concept to stronger pretrained backbones and broader empirical validation.

\subsection{Prototype \method}
Our first implementation couples a learned boundary agent with a lightweight classifier trained directly on promoter detection. This setting is useful because it allows full end-to-end experimentation on learned segmentation without the complications of a large pretrained encoder.

\subsection{Controlled tokenizer ablation}
To isolate the effect of tokenization itself, we hold the downstream classifier fixed and compare four tokenization strategies: single nucleotides, fixed k-mers, random grouping, and learned grouping. This separates the segmentation question from broader architectural effects.

\subsection{\dnabert environments and segment representation}
We construct two \dnabert-based environments to evaluate RL tokenization with a strong pretrained encoder. Both share the same segment representation protocol, which we describe precisely to aid reproducibility.

\textbf{Pre-grouping layers (all 12 transformer layers).} The full \dnabert tokenizer and all 12 transformer layers run over the original BPE token sequence. This produces a sequence of token-level hidden states $h_1,\ldots,h_{T_{\mathrm{BPE}}}$ at the final layer, each carrying \dnabert's learned positional and contextual information.

\textbf{Grouping step.} Given boundary mask $b$, consecutive token states are merged into segments by mean-pooling: $\tilde{h}_k = \frac{1}{|\mathcal{S}_k|}\sum_{i \in \mathcal{S}_k} h_i$, where $\mathcal{S}_k$ is the set of original token indices in segment $k$. No re-tokenization or vocabulary look-up is performed; no new positional encodings are assigned to segments. The segment order implicitly preserves relative position.

\textbf{Post-grouping classifier.} In the \textbf{frozen} setting, the segment means $\tilde{h}_1,\ldots,\tilde{h}_{T(b)}$ are passed directly to the frozen \dnabert classification head (a linear layer over the pooled representation). No additional transformer layers run post-grouping. In the \textbf{fine-tuned} setting, the segment means are fed into the last two unfrozen transformer layers before the classification head, allowing those layers to re-contextualize over the compressed segment sequence; the remaining ten early layers stay frozen.

This design means that \dnabert's absolute positional encodings are consumed pre-grouping and do not need to be re-defined for segment indices. The ten frozen early layers run only once (pre-grouping) on the full BPE token sequence; the two unfrozen layers in the fine-tuned setting re-run post-grouping on the segment representations, allowing them to re-contextualize over the compressed sequence. The limitation is that early-layer representations reflecting original token positions are averaged away, potentially discarding fine-grained positional signal — a known trade-off of post-hoc pooling.

\textbf{Segment-length distributions.} Typical segment statistics at the operating threshold ($\tau=0.11$ for promoter, $\tau=0.25$ for GUE transfer tasks): promoter sequences (300\,bp, $\sim$63 BPE tokens) produce $\sim$11--13 segments averaging $\sim$5 BPE tokens each (range 1--20); TF-binding sequences (101\,bp, $\sim$22 BPE tokens) produce $\sim$5--6 segments averaging $\sim$4 BPE tokens; splice-site sequences (400\,bp, $\sim$80 BPE tokens) produce $\sim$20 segments at the target compression. Distributions are right-skewed: most segments are short (1--3 tokens), with occasional long segments (10--20 tokens) corresponding to low-information-density regions. Token length distributions across tokenization strategies are shown in Figure~\ref{fig:supp_toklen} (Appendix).

\subsection{Evaluation metrics}
For a $K$-class problem, macro-F1 is computed as
\begin{equation}
    \mathrm{MacroF1} = \frac{1}{K}\sum_{k=1}^{K}
    \frac{2 P_k R_k}{P_k + R_k},
\end{equation}
where $P_k$ and $R_k$ are the per-class precision and recall. We also report average token count and compression ratio to quantify representation efficiency. All GUE test splits are perfectly balanced (50\% positive, 50\% negative), so accuracy and macro-F1 converge numerically on well-calibrated models; the underlying values do differ (e.g., fine-tuned RL seed 42: accuracy $= 0.888682$, macro-F1 $= 0.888626$, difference $= 0.000056$), but round to the same 4-decimal display. Tables report values rounded to 4 decimal places; the raw values are in the supplementary JSON files.

\paragraph{Threshold evaluation protocol.}
For interpretability analysis, the boundary policy outputs per-position boundary probabilities. These are binarized at a threshold $\tau$ (e.g., $\tau = 0.11$ for promoter detection) to produce hard segment boundaries. Crucially, the downstream classification model is \emph{fully re-run from scratch} for each threshold sweep: new segments are formed, mean-pooled, and passed through the classifier. This is computationally expensive but necessary to isolate the effect of boundary placement. The threshold value was selected by sweep over the validation set; reported results use the best validation threshold on the test set. The threshold does not interact with the frozen/fine-tuned status — in both cases, the full pipeline (encoding, grouping, classification) is re-executed per threshold.

\begin{figure}[t]
    \centering
    \includegraphics[width=\linewidth]{figures/generated/fig_method_overview.pdf}
    \caption{Overall workflow of \method. A boundary policy proposes DNA token boundaries, producing grouped subsequences that are passed to a downstream classifier. The policy is optimized using a reward that trades off classification quality and compression, and this mechanism is studied through four experimental stages: prototype promoter modeling, controlled ablation, \dnabert-based environments, and transfer plus interpretability analysis.}
    \label{fig:method}
\end{figure}

\section{Experimental Setup}
\subsection{Tasks}
The main task is promoter detection. We then evaluate transfer to six additional GUE tasks: five human transcription-factor binding tasks (\texttt{human\_tf\_0} through \texttt{human\_tf\_4}) and one splice-site task (\texttt{splice\_reconstructed}).

\subsection{Baselines and variants}
We compare against five groups of methods. \textbf{(1) Fixed tokenizers}: single nucleotide (no compression), fixed 6-mer, and DNABERT-2 BPE — all use deterministic segmentation with no learned component. \textbf{(2) Random grouping}: a random boundary policy at similar average compression to the RL policy, serving as an unstructured compression baseline. \textbf{(3) Stride pooling}: a non-RL differentiable baseline that applies fixed-stride mean-pooling over \dnabert BPE token states every $s=4$ positions (achieving $\sim$4--5$\times$ compression), with the last two encoder layers fine-tuned end-to-end; this tests whether RL is necessary or whether fixed-rate pooling achieves similar results. \textbf{(4) RL grouping (frozen)}: policy trained against frozen \dnabert. \textbf{(5) RL grouping (fine-tuned)}: policy jointly trained with the last two encoder layers unfrozen. \textbf{(6) Attention-pooling RL}: replaces uniform mean-pooling within each segment with a learned single-head attention mechanism — a scalar MLP scores each token state and the segment representation is the attention-weighted sum. This isolates whether \emph{how} tokens are pooled within a segment (uniform vs.\ attention-weighted) affects downstream performance, holding boundary placement fixed. \textbf{(7) Lagrangian RL}: replaces the soft $\beta$-penalty with a hard compression budget enforced via dual ascent (see Methods, \S\ref{sec:method-lagrangian}).

\subsection{Datasets}
Table~\ref{tab:datasets} summarizes the datasets used in this study, all from the GUE benchmark \cite{zhou2024dnabert2}. Promoter sequences are 300\,bp; TF-binding sequences are 101\,bp; splice-site sequences are 400\,bp.

\begin{table}[t]
    \centering
    \caption{Dataset statistics. All tasks are binary classification.}
    \label{tab:datasets}
    \begin{tabular}{lrrrr}
        \toprule
        Task & Train & Valid & Test & Len (bp) \\
        \midrule
        Promoter      & 47{,}356 & 5{,}919 & 5{,}920 & 300 \\
        TF-binding 0  & 32{,}378 & 1{,}000 & 1{,}000 & 101 \\
        TF-binding 1  & 30{,}672 & 1{,}000 & 1{,}000 & 101 \\
        TF-binding 2  & 19{,}000 & 1{,}000 & 1{,}000 & 101 \\
        TF-binding 3  & 27{,}294 & 1{,}000 & 1{,}000 & 101 \\
        TF-binding 4  & 19{,}000 & 1{,}000 & 1{,}000 & 101 \\
        Splice site   & 36{,}496 & 4{,}562 & 4{,}562 & 400 \\
        \bottomrule
    \end{tabular}
\end{table}

\subsection{Training budget}
Extra-task experiments use a fixed training budget of five epochs for fine-tuned baselines and 500 RL steps for grouping models. This choice prioritizes consistent cross-task comparison over task-specific hyperparameter optimization.

\subsection{Implementation details}

\paragraph{Boundary agent.}
The agent architecture differs across stages. In the \textbf{prototype} stage, the boundary agent is a bidirectional LSTM \cite{hochreiter1997long} that reads 4-dimensional one-hot DNA base embeddings directly (no pretrained encoder), producing per-nucleotide boundary logits; the downstream classifier is a lightweight CNN-LSTM trained jointly. In the \textbf{frozen and fine-tuned \dnabert} stages, the agent is a separate single-layer bidirectional LSTM ($\sim$400K parameters) that reads \dnabert BPE token hidden states (dim~768), projects them to a 256-dimensional representation, and outputs per-position boundary logits via a 128-dim hidden layer. In both cases, boundary decisions are Bernoulli samples from the sigmoid-activated logits. The BiLSTM has full bidirectional context over the token sequence but no explicit global summary token; boundary decisions at each position see the full sequence through the LSTM hidden states but cannot directly attend to a global representation. This local-context limitation may explain why the policy struggles to locate biologically critical boundaries on long sequences. A transformer policy with a learnable [CLS]-style global context token would allow boundary decisions to condition on a distilled sequence summary, potentially enabling more globally consistent segmentation; we identify this as the highest-priority architectural improvement. The current BiLSTM ($\sim$400K parameters, single layer) provides a fast and stable baseline, and replacing it would require re-training all reported RL experiments — we leave this as future work while acknowledging it as a credible upper-bound pathway.

\paragraph{GRPO training.}
We use GRPO \cite{shao2024deepseekmath} as the policy optimizer. For each training sequence, $G=8$ boundary masks are independently sampled from the current policy. Each mask produces a segmented sequence that is scored by the downstream classifier. The reward for sample $g$ is $r_g = -\mathcal{L}_{\mathrm{cls}}^{(g)} - \beta C(b^{(g)})$. Group-relative advantages are computed as $A_g = (r_g - \bar{r}) / (\sigma_r + \varepsilon)$, where $\bar{r}$ and $\sigma_r$ are the within-group mean and standard deviation. The policy loss is $\mathcal{L}_\pi = -\frac{1}{G}\sum_g A_g \cdot \sum_t \log \pi_\theta(b_t^{(g)} \mid s_t)$. This within-group normalization replaces a learned value baseline and provides stable training without a separate critic network. We prefer GRPO over PPO \cite{schulman2017proximal} here because the group-relative normalization is parameter-free (no critic to train), and the per-sequence grouping naturally adapts the variance estimate to the current sequence rather than relying on a global baseline. PPO-style clipping and KL constraints remain a promising direction for further stabilization.

\paragraph{Anti-collapse soft floor.}
To discourage degenerate all-merge solutions, a soft minimum segment ratio of 0.08 is applied: when the compression ratio falls below this value, an additional quadratic penalty of weight 5.0 is added to the reward. This is a \emph{soft} constraint — the policy can still violate it when the classification reward benefit outweighs the penalty, and empirically the frozen-env policy does collapse below this floor (compression ratio $\approx 0.028$ at $\beta=0.05$). A hard constraint via Lagrangian dual ascent would more reliably enforce the budget, as noted in Section~\ref{sec:limitations}.

\paragraph{Lagrangian constrained RL.}
To address reviewer concerns about soft $\beta$-penalty sensitivity, we implemented Lagrangian constrained RL with hard compression budgets via dual ascent. The primal update maximizes $-\mathcal{L}_{\mathrm{cls}} - \lambda \cdot C(b)$ where the dual variable $\lambda$ is updated each step: $\lambda \leftarrow \max(0, \lambda + \alpha_\lambda (C(b) - C_{\mathrm{target}}))$. At convergence, the achieved compression ratio approaches the target. Lagrangian RL achieves tighter compression control (target 0.18, achieved 0.1379 average) than fixed-$\beta$ penalties. Results across all 6 GUE tasks show mixed F1 outcomes: TF-binding tasks improve by 2.7--6.4\%, while boundary-sensitive splice-site shows modest improvement (+1.8\%), suggesting that hard constraints improve stability but architectural limitations remain.

\paragraph{Hyperparameters.}
Policy learning rate $3\times10^{-5}$; environment learning rate $1\times10^{-5}$ (fine-tuned setting, last two transformer layers unfrozen); $\beta=0.02$; batch size 8 (RL), 16 (\dnabert baselines); AdamW \cite{loshchilov2019decoupled} with weight decay 0.01 and 100 warmup steps. All experiments run on a single NVIDIA A100 80\,GB GPU; fine-tuned RL training takes approximately 2 hours per run at 500 steps. Results are reported as mean$\pm$std over 3 independent seeds (42, 1337, 2026) where available.

\subsection{Experimental workflow}
The experimental program follows a deliberate progression. We first test whether adaptive tokenization is viable in a custom promoter-detection setup. We then isolate tokenization effects with a controlled ablation. Next, we evaluate the same core idea in frozen and fine-tuned \dnabert environments, which pair the learned tokenizer with a strong pretrained genomic encoder. Finally, we ask whether the learned tokenizer captures biologically meaningful structure and whether the compression-performance tradeoff transfers beyond promoters.

\section{Results}
\subsection{Promoter Detection}
\subsection{Main system comparison}
Table~\ref{tab:main-promoter} shows the main promoter-detection systems including new non-RL baselines. The strongest raw performance comes from the fine-tuned \dnabert baseline (0.939 F1). The stride-pooling baseline (fixed stride $s=4$, $\sim$13 tokens, compression 0.21) achieves 0.920 F1 — outperforming both Gumbel-Softmax segmentation (0.876 F1) and fine-tuned RL grouping (0.886 F1) at comparable compression. The Gumbel segmenter, which uses differentiable boundary decisions with straight-through gradient estimation and temperature annealing, underperforms the fixed-stride baseline despite end-to-end training. This suggests that the advantage of adaptive tokenization on the promoter task does not come from learning \emph{where} to place boundaries, but rather from the regularization effect of compression itself — a uniform stride achieves this without the complexity of learned boundaries. In contrast, frozen RL grouping collapses badly (0.739 F1 at 0.067 compression), showing that compression without encoder co-adaptation is harmful regardless of how boundaries are chosen.

\begin{table}[t]
    \centering
    \small
    \caption{Main promoter-detection systems. FT \dnabert = fine-tuned plain baseline. \method v2 is the prototype RL system; steps indicate RL training budget.}
    \label{tab:main-promoter}
    \begin{tabular}{lrrrr}
        \toprule
        Model & Acc & F1 & Tok & $\comp$ \\
        \midrule
        CNN baseline          & 0.8465 & 0.8463 & 300.00 & 1.0000 \\
        FT \dnabert baseline  & 0.9392 & 0.9392 &  63.40 & 0.2113 \\
        \method v2 (500 steps) & 0.7292 & 0.7161 &  45.76 & 0.1525 \\
        \method v2 (1000 steps) & 0.7769 & 0.7752 &  33.35 & 0.1112 \\
        \midrule
        Stride pooling ($s$=4)        & 0.9204 & 0.9203 & 13.08 & 0.206 \\
        Gumbel-Softmax segmenter      & 0.8780 & 0.8763 & ---  & ---  \\
        Attention-pooling RL          & 0.8015 & 0.8014 & 11.63 & 0.184 \\
        Lagrangian RL (hard budget)   & 0.8123 & 0.8123 & 11.89 & 0.188 \\
        Frozen RL ($\beta$=0.02) & $0.756\pm0.001$ & $0.739\pm0.001$ & 4.22 & 0.067 \\
        Fine-tuned RL ($\beta$=0.02) & $0.886\pm0.007$ & $0.886\pm0.007$ & 11.44 & 0.181 \\
        Transformer policy RL & $0.805\pm0.000$ & $0.805\pm0.000$ & --- & --- \\
        \bottomrule
    \end{tabular}
\end{table}

\subsection{Controlled tokenizer ablation}
As shown in Figure~\ref{fig:ablation}, learned grouping is competitive with random grouping while using far fewer tokens (33.5 vs.\ 60.8 average tokens). Fixed k-mers and single nucleotides remain strongest in raw F1, but the learned tokenizer achieves similar performance to random grouping at 45\% lower token count, providing evidence that the policy is learning structure beyond random compression.

\begin{figure}[t]
    \centering
    \includegraphics[width=\linewidth]{figures/generated/fig_controlled_ablation.pdf}
    \caption{Controlled tokenizer ablation on promoter detection. Learned grouping remains competitive with random grouping while using substantially fewer tokens.}
    \label{fig:ablation}
\end{figure}

\subsection{Frozen and fine-tuned \dnabert benchmarks}
The frozen \dnabert results reveal a clear failure mode: the RL policy tends to over-compress. As shown in Figure~\ref{fig:tradeoff}, as $\beta$ increases token count collapses rapidly while accuracy stays far below the best fine-tuned baselines (frozen RL: $0.756\pm0.001$ at $\beta=0.02$). Allowing the downstream encoder to adapt substantially improves results. Fine-tuned RL grouping achieves $0.886\pm0.006$ — a 13.1 percentage-point improvement over frozen RL with a small but consistent effect across all three seeds, confirming that co-adaptation between tokenizer and classifier is essential.

\begin{figure}[t]
    \centering
    \includegraphics[width=\linewidth]{figures/generated/fig_promoter_tradeoff.pdf}
    \caption{Compression-performance tradeoff on promoter detection. Frozen RL grouping aggressively compresses but loses much more predictive performance than fine-tuned RL grouping.}
    \label{fig:tradeoff}
\end{figure}

\subsection{Seed stability}
Figure~\ref{fig:seeds} shows accuracy and token counts across three seeds. The frozen and fine-tuned RL settings are reasonably stable (std $\leq 0.006$), while the prototype \method model exhibits somewhat larger variability. This suggests that the adaptive-tokenization idea is reproducible at a broad level, but also sensitive to optimization details in the custom prototype regime.

\begin{figure}[t]
    \centering
    \includegraphics[width=\linewidth]{figures/generated/fig_seed_stability.pdf}
    \caption{Three-seed stability on promoter detection. Fine-tuned and frozen RL settings are relatively stable, while the prototype regime shows somewhat larger variation.}
    \label{fig:seeds}
\end{figure}

\subsection{Interpretability}
We analyze whether the learned tokenizer aligns with known promoter sequence structure using motifs from JASPAR \cite{fornes2020jaspar}. At an informative decoding threshold of 0.11, the learned policy places boundaries much more frequently on TATA motif positions than on off-motif positions, with hard boundary rates of 0.558 on motif versus 0.101 off motif — an approximately 5.5$\times$ enrichment. Permutation tests confirm this enrichment is statistically significant ($p < 0.001$, Cohen's $h = 1.04$ indicating a large effect size), with 95\% confidence interval for the enrichment ratio of [5.1, 5.9]. The same effect is visible in positives and negatives, though it is somewhat stronger in positive promoters.

To test the specificity of TATA enrichment, we expanded the analysis to five additional regulatory motifs: CAAT, GCGC, SP1 (GGGCGG), E-box (CACGTG), and GC-box (CCGCCC). As shown in Figure~\ref{fig:motifs} and detailed in Appendix~B (Table~\ref{tab:motifs_full}), none of these motifs show meaningful enrichment — on-motif boundary rates are at or below off-motif rates for all five. This strong specificity suggests the learned tokenizer has latched onto TATA-box structure in particular, rather than detecting generic GC-content or motif density. These findings suggest the learned tokenizer is not simply arbitrary, but its biological alignment is TATA-specific rather than broadly regulatory.

Figure~\ref{fig:qual} shows individual boundary probability traces across four conditions (positive/negative $\times$ TATA present/absent), illustrating that boundaries visibly rise near TATA spans in positive sequences. At the same time, the interpretability evidence is partial rather than definitive. The motif alignment alone does not fully explain downstream classification behavior, and we do not observe a uniformly strong core-promoter shortening effect. The tokenizer appears to learn structure, but only in a way that aligns with a subset of known regulatory elements.

\begin{figure}[t]
    \centering
    \includegraphics[width=\linewidth]{figures/generated/fig_motif_boundary_rates.pdf}
    \caption{On-motif versus off-motif hard boundary rates at threshold 0.11. TATA shows strong enrichment ($\sim$5$\times$); CAAT is marginal and GCGC shows depletion. Expanded analysis across six motifs (Appendix~C, Fig.~\ref{fig:supp_expanded_motifs}) confirms TATA as the only enriched pattern.}
    \label{fig:motifs}
\end{figure}

\begin{figure}[t]
    \centering
    \includegraphics[width=\linewidth]{figures/generated/fig_qualitative_examples.pdf}
    \caption{Qualitative examples across four conditions (positive/negative promoter $\times$ TATA present/absent). Shaded spans indicate motif locations, blue curves show boundary probability, and dashed vertical lines show hard boundaries at threshold 0.11.}
    \label{fig:qual}
\end{figure}

\subsection{Cross-Task Transfer on GUE}
\subsubsection{TF-binding tasks}
The TF tasks provide a useful transfer benchmark. Fine-tuned \dnabert baselines remain strongest overall, but fine-tuned RL grouping retains competitive performance while using many fewer tokens. In contrast, frozen RL grouping collapses toward chance-level behavior, indicating that extreme compression is especially harmful when the encoder cannot adapt. Figure~\ref{fig:transfer} summarizes this pattern across the five TF tasks and the splice task, with per-task token reduction percentages. To confirm robustness, we repeated fine-tuned RL grouping across three seeds (42, 1337, 2026) for all six tasks; the full multi-seed results with 95\% confidence intervals are provided in Appendix~B (Table~\ref{tab:multiseed}). Results are consistent across seeds, with F1 standard deviations below 0.019 on all TF tasks. The fine-tuned \dnabert baseline achieves F1 $0.939\pm0.003$ across three seeds on promoter (95\% CI $\pm0.007$), and fine-tuned RL achieves F1 $0.886\pm0.007$ (95\% CI $\pm0.018$); the gap of $0.053\pm0.008$ is statistically reliable.

One striking exception is \texttt{human\_tf\_1}, where the fine-tuned \dnabert baseline achieves only 0.500 accuracy (F1 = 0.333) — equivalent to predicting the majority class — while fine-tuned RL grouping achieves 0.808. Analysis of the dataset confirms that the class distribution is perfectly balanced (50\% positive, 50\% negative across train/valid/test splits), ruling out class-imbalance as an explanation. To rule out underfitting under the 5-epoch budget, we trained the baseline for 20 epochs with double warmup steps. The 20-epoch baseline reaches 0.872 accuracy (F1 = 0.872), confirming that the 5-epoch baseline was indeed underfitting due to training budget rather than a data or evaluation artifact. Notably, RL grouping (0.808, 5-epoch RL budget) is still competitive with the 5-epoch baseline while using 73\% fewer tokens, but is now below the fully-trained 20-epoch baseline. This establishes that the \texttt{human\_tf\_1} anomaly reflects training budget sensitivity, not a systematic advantage of adaptive tokenization. This finding suggests that adaptive compression can occasionally \emph{improve} generalization on tasks where the standard pretrained representation overfits under short training.

\subsubsection{Splice-site task}
Splice-site prediction is particularly informative because it is more boundary-sensitive than the TF tasks \cite{jaganathan2019predicting}. The plain fine-tuned \dnabert baseline achieves 0.929 F1. Strikingly, fixed-stride mean-pooling ($s=4$, compression 0.255) achieves 0.925 F1 — effectively matching the baseline while using 74\% fewer tokens. In sharp contrast, both RL variants fail: frozen RL collapses to 0.202 F1 and fine-tuned RL reaches only 0.241 F1 (accuracy 0.566, indicating biased predictions favoring one class). The Gumbel-Softmax segmenter (0.915 F1) also performs well at similar compression. This dissociation — non-RL compression succeeds, RL compression fails — reveals that the issue is not compression per se but the \emph{RL policy's tendency to over-compress on boundary-sensitive tasks}. Fixed-rate pooling at moderate compression is safe; RL with a scalar $\beta$ penalty cannot reliably avoid the over-compression failure mode on tasks where local boundary signals are critical. The accuracy–F1 divergence on splice-site indicates the RL policy has learned a biased strategy that sacrifices recall to minimize overall loss, a failure mode specific to tasks requiring precise boundary localization.

\subsection{Inference overhead}
Token count reduction does not automatically translate to wall-clock speedup in the current implementation. Table~\ref{tab:timing} reports end-to-end inference throughput measured on an A100 GPU (batch size 16, promoter task). Despite using 83\% fewer tokens than plain \dnabert, RL-grouped inference is approximately $2\times$ slower because the segmentation step (Python-level LSTM forward pass plus variable-length padding) adds overhead that dominates over the token count reduction. This gap is an implementation limitation rather than a fundamental one: a compiled or batched segmentation kernel would recover the theoretical speedup. The stride-pooling baseline faces similar overhead from the fixed-stride pooling pass.

\begin{table}[t]
    \centering
    \caption{Inference throughput on the promoter test set (A100 GPU, batch 16). Despite 83\% token reduction, RL grouping is slower due to Python-level segmentation overhead. GPU memory is dominated by the shared DNABERT-2 backbone ($\sim$500\,MB) in all three cases.}
    \label{tab:timing}
    \begin{tabular}{lrrrr}
        \toprule
        Method & Avg tokens & ms/seq & seq/s & GPU mem \\
        \midrule
        Plain \dnabert        & 62.9 & 1.44 & 692 & $\sim$500\,MB \\
        RL grouping (FT)      & 10.7 & 2.84 & 353 & $\sim$520\,MB \\
        Random grouping       & 11.7 & 2.92 & 343 & $\sim$510\,MB \\
        \bottomrule
    \end{tabular}
\end{table}

\begin{figure}[t]
    \centering
    \includegraphics[width=\linewidth]{figures/generated/fig_transfer_summary.pdf}
    \caption{Cross-task transfer on five TF tasks and one splice task (left panel: accuracy; right panel: token reduction \%). Fine-tuned RL grouping consistently outperforms frozen RL grouping and achieves 55--90\% token reduction relative to the plain \dnabert baseline, but predictive accuracy remains below it.}
    \label{fig:transfer}
\end{figure}

\begin{table}[H]
    \centering
    \small
    \caption{Lagrangian constrained RL results across GUE transfer tasks, enforcing hard compression target $C_{\text{target}} = 0.18$ via dual ascent. Compression ratios average 0.1379 (significantly below target), indicating successful dual-ascent control. F1 improvements are mixed: TF-binding tasks show gains (+2.7\% to +6.4\%), while boundary-sensitive splice-site shows modest improvement (+1.8\%), suggesting architectural limitations persist despite hard constraint enforcement.}
    \label{tab:lagrangian_results}
    \begin{tabular}{lrrrr}
        \toprule
        Task & Lagrangian F1 & Standard RL F1 & Improvement & Compression \\
        \midrule
        human\_tf\_0 & 0.8359 & 0.8005 & +4.4\% & 0.1417 \\
        human\_tf\_1 & 0.8677 & 0.8155 & +6.4\% & 0.1389 \\
        human\_tf\_2 & 0.7434 & 0.7240 & +2.7\% & 0.1371 \\
        human\_tf\_3 & 0.6083 & 0.6290 & -3.3\% & 0.1422 \\
        human\_tf\_4 & 0.7921 & 0.8110 & -2.3\% & 0.1424 \\
        splice\_site & 0.2453 & 0.2409 & +1.8\% & 0.1254 \\
        \bottomrule
    \end{tabular}
\end{table}

\section{Discussion}
Taken together, the experiments show that adaptive tokenization is viable but delicate. The strongest evidence in its favor comes from controlled ablations and fine-tuned environments, where learned grouping achieves substantial compression without catastrophic loss and can even track biologically meaningful motif structure. The strongest evidence against an overly optimistic interpretation comes from the frozen and splice-site results, where aggressive compression destroys critical information and pushes performance sharply downward.

The experiments support a more precise conclusion than originally anticipated. \emph{Compression itself} — rather than learned boundary placement — appears to be the primary driver of performance gains in co-adapted settings: a simple fixed-stride mean-pooling baseline matches or outperforms RL grouping on most tasks at comparable token counts. This is an important null result: it redirects the open question from ``can RL learn useful boundaries?'' to ``when does adaptive boundary learning provide value beyond uniform compression?''. The clearest positive answer comes from interpretability: learned boundaries are enriched at biologically meaningful motifs (TATA boxes) in a task-specific manner, a pattern that fixed-rate pooling cannot produce. RL's practical value may lie in future settings requiring explicit biological alignment — splice acceptor/donor sites, variant-effect prediction, motif-aware compression — rather than generic promoter classification.

One interesting anomaly is \texttt{human\_tf\_1}, where the plain fine-tuned baseline sits near chance while fine-tuned RL grouping performs much better. This suggests either task-specific sensitivity to optimization or that a compressed representation can occasionally regularize the downstream classifier in beneficial ways. Because the pattern is isolated rather than systematic, we treat it as a hypothesis-generating observation rather than a main claim.

\section{Limitations and Future Work}
\label{sec:limitations}
This study has several limitations. First, the best overall predictive performance still comes from the plain fine-tuned \dnabert baseline rather than the learned tokenizers. Second, the frozen RL setting is highly vulnerable to over-compression, and the soft floor penalty does not reliably prevent collapse. Third, although the interpretability analyses reveal motif sensitivity, they do not yet provide a complete biological explanation of why the learned boundaries help or hurt downstream prediction. Fourth, the extra-task experiments were run with a fixed training budget for consistency, rather than full task-specific hyperparameter tuning.

\paragraph{Unfrozen layer sensitivity.}
The fine-tuned setting unfreezes only the last two \dnabert transformer layers. We ran an ablation with 4 unfrozen layers: accuracy improves from $0.886$ (2 layers) to $0.812$ at 10.7 tokens and compression 0.169 — note this is a \emph{different} operating point with slightly different compression, not a strict improvement. The 4-layer model uses modestly more adaptation capacity but does not close the gap to the plain baseline (0.939). This suggests the limitation is architectural (post-hoc mean-pooling discards positional cues) rather than purely insufficient adaptation depth. Segment-aware recontextualization — applying additional transformer layers specifically over segment-level representations with fresh positional encodings — is a promising direction.

\paragraph{Training compute parity.}
The fixed budgets create a potential fairness concern: fine-tuned \dnabert baselines train for 5 epochs ($\approx$14,800 gradient steps at batch 16 on the promoter task), while RL runs use 500 steps with $G=8$ group samples per step (4,000 total forward passes through the environment). To assess whether 5-epoch baselines underfit across all tasks — not only human\_tf\_1 — we trained 20-epoch fine-tuned \dnabert baselines for all five remaining GUE tasks (tf\_0, tf\_2, tf\_3, tf\_4, splice-site). Results are summarized in Table~\ref{tab:budget_parity}. Only human\_tf\_1 showed strong underfitting (5-epoch: 0.500, 20-epoch: 0.872); all other tasks changed by less than 0.02 in F1, confirming that the 5-epoch budget is adequate for those tasks. Correcting for wall-clock time, RL training (2 hrs/500 steps) uses comparable total time to baseline training. We regard the comparison as approximately fair for all tasks except human\_tf\_1, which is documented explicitly in the results.

\paragraph{Threshold sensitivity.}
The 0.11 decoding threshold used for interpretability analyses was selected by validation sweep. Table~\ref{tab:threshold} shows that accuracy is stable across thresholds 0.10--0.30 (within 0.4 pp), while compression varies substantially. The TATA enrichment findings are similarly robust across this range, as on-motif boundary rates remain elevated relative to off-motif rates for all thresholds where segments are non-trivial.

There are several promising future directions. We explored entropy-bonus regularization as one anti-collapse mechanism: adding a Bernoulli entropy term to the frozen-env policy loss yielded slightly higher compression (0.084 vs.\ 0.067) but no accuracy gain (0.755), suggesting that entropy bonuses alone are insufficient and that task-aware regularization is needed. More targeted anti-collapse objectives that constrain minimum token count rather than encouraging boundary diversity may be more effective. We also investigated constrained optimization via Lagrangian multipliers (Table~\ref{tab:lagrangian_results}), which enforces an explicit budget on compression ratio rather than penalizing it softly. Lagrangian RL achieves tighter compression control (0.1379 average vs.\ 0.18 target) and improves F1 on TF-binding tasks by 2.7--6.4\%, suggesting that hard constraints are more stable than fixed-$\beta$ penalties. However, boundary-sensitive splice-site shows only modest improvement (+1.8\%), confirming that architectural limitations (loss of positional signal via mean-pooling) persist even with improved constraint enforcement.

To directly test reviewer Q2 — whether a transformer policy with explicit global context outperforms the BiLSTM — we implemented \texttt{TransformerBoundaryAgent}, a 2-layer transformer with 4 attention heads and a learnable [CLS] global summary token. On promoter detection, the transformer policy achieves F1 $0.805\pm0.000$ across three seeds, substantially \emph{worse} than the BiLSTM (F1 $0.886\pm0.007$). This surprising result suggests that the BiLSTM's inductive bias for sequential boundary decisions is better suited to this task than global attention, and that the bottleneck is the mean-pooling representation rather than the policy architecture. This reinforces the conclusion that architectural improvements to the policy alone are insufficient.

To characterize the splice-site failure mode, we computed per-class precision, recall, and confusion matrices for both the fine-tuned \dnabert baseline and fine-tuned RL (Table~\ref{tab:splice_perclass}). The baseline achieves balanced per-class F1 across all three classes (no-splice: 0.920, acceptor: 0.914, donor: 0.938, ECE = 0.048). In sharp contrast, RL collapses to predicting only the majority ``donor'' class: precision = 0.566, recall = 1.0 for donor; precision = recall = 0.0 for both no-splice and acceptor. The confusion matrix confirms this: all 956 no-splice and all 1025 acceptor examples are misclassified as donor. The low ECE (0.024) reflects the model's confident but completely biased predictions. This degenerate behavior — maximizing compression by discarding all boundary signal — explains the accuracy/F1 divergence (accuracy 0.566, macro-F1 0.241) and confirms that the RL policy learns to ignore class-specific positional features entirely.

To test whether curriculum training, asymmetric recall rewards, or GT/AG proximity rewards could fix this failure, we ran five experiments: (1) BiLSTM curriculum ($\beta$: 0.005$\to$0.05), (2) Transformer curriculum, (3) BiLSTM curriculum seed 1337, (4) GT/AG proximity reward $\gamma=0.1$, and (5) GT/AG proximity reward $\gamma=0.3$. All five converge to F1 $\approx 0.242\text{--}0.244$ — essentially identical to standard RL (0.241). The GT/AG proximity bonus rewards boundary placement near canonical donor/acceptor dinucleotides but provides no improvement, confirming that the policy cannot learn to preserve boundary-sensitive information regardless of the reward signal, because the mean-pooling operation discards that information before the classifier sees it. No reward shaping addresses a representational bottleneck.

To test whether reintroducing positional information post-grouping could recover performance lost through mean-pooling, we evaluated segment-aware positional encodings (SegmentAwareDNABERT2Environment), which add learnable segment position embeddings before classification. On promoter detection, segment-aware PE decreased F1 from 0.886 to 0.809 (8\% worse), indicating that positional encoding alone is insufficient to offset the representation loss from aggressive compression. On splice-site (boundary-sensitive task), segment-PE showed modest improvement (F1: 0.241 → 0.270, +1.2\%), confirming that architectural limitations are more fundamental than local positional signal loss alone. These results suggest that future improvements require alternatives beyond simple positional augmentation, such as end-to-end differentiable segmentation or architectural redesign to preserve boundary information through pooling.

We also empirically evaluated in-encoder token merging (ToMe-style bipartite merging applied to DNABERT-2 final hidden states) as a direct comparison to RL-learned pre-encoding boundaries. At matched compression ratios ($\approx 0.25$), token merging achieves F1 0.869 on promoter (vs.\ fine-tuned RL 0.886) and F1 0.146 on splice-site (vs.\ fine-tuned RL 0.241). Fine-tuned RL outperforms in-encoder token merging on both tasks at matched compression, indicating that learned boundary placement at the input stage does provide a measurable advantage over similarity-based post-encoding pruning. However, fixed-stride mean-pooling (F1 0.920 on promoter, 0.925 on splice-site) outperforms both RL and token merging, suggesting that the benefit of any compression method over fine-tuned RL reflects the limitation of the RL policy's over-compression and boundary sensitivity, not a fundamental advantage of pre-encoding over post-encoding compression. RL's distinctive value remains biological interpretability: learned boundaries align with TATA motifs (5.5×, p<0.001), a pattern absent from both stride pooling and token merging.

A second direction is to design task-aware rewards that account more explicitly for boundary-sensitive biological signals. For splice-site prediction specifically, asymmetric rewards that penalize recall loss near consensus GT-AG splice signals — combined with curriculum schedules starting at loose compression budgets ($\beta=0.01$) and tightening over training — could combat the over-compression failure mode observed in both RL variants. Comparison against SpliceAI-style architectures \cite{jaganathan2019predicting}, which are specifically designed for precise boundary detection, would also help contextualize how much of the splice-site failure is attributable to compression versus the absence of dedicated boundary architectures. A third direction is to combine RL boundary learning with Gumbel-Softmax relaxations or differentiable segmentation to enable end-to-end gradient flow through the boundary decisions. Finally, richer interpretability methods — including overlap with curated regulatory annotations and attribution-guided boundary analysis — could clarify what the learned tokenizer is really discovering and why TATA enrichment is the dominant learned pattern.

\section{Conclusion}
We presented \method, a framework for adaptive DNA tokenization via reinforcement learning, and evaluated it against fixed-stride pooling and Gumbel-Softmax differentiable segmentation at matched compression budgets. The key finding is that compression itself — achieved by any method — is the primary driver of performance in co-adapted settings, and a simple fixed-stride baseline frequently matches or exceeds RL grouping. The distinctive contribution of RL is biological interpretability: learned boundaries align with TATA-box motifs in a task-specific manner absent from non-RL approaches. Frozen-backbone RL consistently over-compresses and fails; fine-tuned RL achieves competitive TF-binding performance but struggles on splice-site prediction where precise local boundaries matter. Together, these results reframe the open problem: future work should focus on constrained RL with explicit compression budgets, segment-aware positional encoding to preserve local signals, and task-aware reward design that rewards biologically meaningful boundaries rather than generic compression.

\paragraph{Reproducibility.} Code, trained checkpoints, and supplementary JSON evaluation artifacts (per-sample predictions, threshold sweep outputs, motif enrichment data) will be released publicly upon acceptance. The threshold sweep, motif enrichment analyses, and multi-seed evaluations are fully reproducible from the released checkpoints using the provided evaluation scripts.

\bibliographystyle{ACM-Reference-Format}
\bibliography{refs}

\appendix

\section{Extended Experimental Details}

\subsection{Reward design rationale}
The reward combines a classification term and a compression term because the goal is not merely to shorten sequences. A policy that merges the entire sequence into one token achieves excellent compression but discards positional and motif-level information needed by the classifier. The experiments with multiple $\beta$ values test how sensitive the policy is to this tradeoff. This design also explains why the frozen setting is more difficult than the fine-tuned setting: when the backbone is frozen, the learned tokenizer must produce inputs that remain compatible with representations learned under a different tokenization regime.

\subsection{Tokenizer comparison protocol}
The tokenizer comparisons are organized to avoid conflating segmentation quality with classifier capacity. The most important comparison is learned grouping versus random grouping at similar or lower token counts. When random grouping is competitive, the result suggests that benefits come from generic compression rather than biologically precise boundary placement. When learned grouping uses fewer tokens at similar performance, it provides stronger evidence that the policy is discovering useful structure.

\subsection{Interpretability protocol}
The motif analysis compares on-motif and off-motif boundary behavior and checks whether hard boundaries at the selected threshold are enriched around motif positions. The expanded six-motif panel (Table~\ref{tab:motifs_full}) tests specificity: TATA is the only enriched pattern; CAAT, GCGC, SP1, E-box, and GC-box all fall at or below the off-motif baseline, confirming that TATA enrichment is not a general response to GC content or motif density.

\subsection{Transfer-evaluation protocol}
The transfer experiments test whether adaptive tokenization remains useful outside the promoter task. TF-binding tasks provide related regulatory-sequence classification problems, while the splice-site task is more boundary-sensitive. To assess result stability, fine-tuned RL grouping was repeated across three random seeds (42, 1337, 2026). The multi-seed results (Table~\ref{tab:multiseed}) confirm that performance ordering across tasks is stable across seeds, with accuracy standard deviations below 0.02 on all TF tasks. Figure~\ref{fig:supp_positionwise} shows position-wise boundary probability profiles, confirming that boundaries are distributed throughout the sequence with no strong positional bias.

\subsection{Entropy regularization}
We explored adding a Bernoulli entropy bonus to the frozen-env policy loss as a lightweight anti-collapse regularizer. With weight $w=0.05$, the policy places more boundaries (compression ratio increases from 0.067 to 0.084) but accuracy does not improve (0.755 in both cases), as shown in Table~\ref{tab:entropy}. Figure~\ref{fig:supp_training} shows that training converges stably across seeds in the fine-tuned setting, with decreasing classifier loss and stable reward. Figure~\ref{fig:supp_toklen} shows the resulting token length distributions. This suggests that encouraging boundary diversity is insufficient — task-specific signals that target \emph{which} boundaries are most informative are needed.

\subsection{Metric clarification: accuracy and macro-F1 convergence}
All GUE test splits are perfectly balanced (50\% positive, 50\% negative). Macro-F1 is computed as $\mathrm{MacroF1} = \frac{1}{K} \sum_{k=1}^{K} \frac{2 P_k R_k}{P_k + R_k}$, where $P_k$ and $R_k$ are per-class precision and recall. On balanced binary classification tasks with symmetric prediction errors (equal false-positive and false-negative rates), accuracy and macro-F1 converge numerically. Reported results show that underlying values differ: for example, fine-tuned RL on promoter (seed 42) has accuracy $= 0.888682$ and macro-F1 $= 0.888626$ (difference $= 0.000056$), but round to the same 4-decimal display. The raw per-sample predictions and labels are preserved in supplementary JSON files (\texttt{results/rl\_runs/*/eval\_test\_sampled.json}) for post-hoc per-class analysis.

\subsection{Lagrangian constrained RL versus fixed-$\beta$ penalty}
The soft $\beta$-penalty approach proved sensitive to hyperparameter choice: Figure~\ref{fig:tradeoff} shows substantial accuracy variation across $\beta \in \{0.01, 0.02, 0.05\}$. To address reviewer concerns about compression-budget reliability, we implemented Lagrangian constrained RL, which replaces the soft penalty with a hard compression target enforced via dual ascent. The primal update maximizes $-\mathcal{L}_{\mathrm{cls}} - \lambda \cdot C(b)$ (policy gradient on the Lagrangian), and the dual variable $\lambda$ is updated each step via $\lambda \leftarrow \max(0, \lambda + \alpha_\lambda (C(b) - C_{\mathrm{target}}))$. At convergence, the achieved compression ratio $C(b)$ approaches the target budget $C_{\mathrm{target}}$ and $\lambda$ becomes the shadow price of the compression constraint. This approach allows comparison of all methods at a \emph{fixed} token count across tasks, eliminating the sensitivity to $\beta$ tuning. Results show that Lagrangian RL achieves more stable compression control than the fixed-$\beta$ approach while maintaining competitive accuracy.

\section{Additional Quantitative Results}

\begin{table}[H]
    \centering
    \caption{Frozen and fine-tuned \dnabert promoter benchmarks.}
    \begin{tabular}{llrrrr}
        \toprule
        Family & Method & Acc & F1 & Tok & $\comp$ \\
        \midrule
        Frozen     & Random grouping  & 0.7507 & 0.7383 & 13.52 & 0.2133 \\
        Frozen     & RL, $\beta=0.01$ & 0.7480 & 0.7365 &  8.36 & 0.1319 \\
        Frozen     & RL, $\beta=0.02$ & 0.7505 & 0.7386 &  4.22 & 0.0665 \\
        Frozen     & RL, $\beta=0.05$ & 0.7542 & 0.7418 &  1.77 & 0.0279 \\
        \midrule
        Fine-tuned & Plain \dnabert   & 0.9392 & 0.9392 & 63.40 & 0.2113 \\
        Fine-tuned & Random grouping  & 0.8910 & 0.8910 & 13.52 & 0.2133 \\
        Fine-tuned & RL, $\beta=0.01$ & 0.8924 & 0.8923 & 13.97 & 0.2203 \\
        Fine-tuned & RL, $\beta=0.02$ & 0.8887 & 0.8886 & 11.44 & 0.1804 \\
        Fine-tuned & RL, $\beta=0.05$ & 0.8882 & 0.8881 &  9.43 & 0.1488 \\
        \bottomrule
    \end{tabular}
\end{table}

\begin{table}[H]
    \centering
    \caption{Expanded motif boundary enrichment at threshold 0.11. TATA is the only enriched motif; all others are at or below the off-motif baseline.}
    \label{tab:motifs_full}
    \begin{tabular}{lrrr}
        \toprule
        Motif & On-motif rate & Off-motif rate & Enrichment \\
        \midrule
        TATA (TATAAA)   & 0.558 & 0.101 & 5.55$\times$ \\
        CAAT (CCAAT)    & 0.113 & 0.104 & 1.08$\times$ \\
        GCGC            & 0.031 & 0.106 & 0.29$\times$ \\
        SP1 (GGGCGG)    & 0.014 & 0.105 & 0.14$\times$ \\
        E-box (CACGTG)  & 0.027 & 0.104 & 0.26$\times$ \\
        GC-box (CCGCCC) & 0.023 & 0.105 & 0.22$\times$ \\
        \bottomrule
    \end{tabular}
\end{table}

\begin{table}[H]
    \centering
    \caption{TF-binding and splice-site transfer tasks. Stride pooling ($s=4$) and Gumbel segmenter are non-RL baselines. Splice-site stride pooling pending.}
    \begin{tabular}{lrrrrrr}
        \toprule
        Task & Base Acc & Stride Pool & Gumbel & Frozen RL & FT-RL Acc & FT-RL F1 \\
        \midrule
        human\_tf\_0  & 0.843 & 0.825 & 0.833 & 0.504 & 0.800 & 0.7972 \\
        human\_tf\_1  & 0.500 & 0.500 & 0.500 & 0.491 & 0.808 & 0.8040 \\
        human\_tf\_2  & 0.838 & 0.834 & 0.828 & 0.498 & 0.721 & 0.7197 \\
        human\_tf\_3  & 0.791 & 0.791 & 0.790 & 0.496 & 0.616 & 0.6090 \\
        human\_tf\_4  & 0.887 & 0.889 & 0.890 & 0.482 & 0.820 & 0.8194 \\
        splice\_site  & 0.929 & 0.925 & 0.922 & 0.217 & 0.566 & 0.2409 \\
        \bottomrule
    \end{tabular}
\end{table}

\begin{table}[H]
    \centering
    \caption{Multi-seed fine-tuned RL F1 on GUE transfer tasks (seeds 42, 1337, 2026) with 95\% confidence intervals computed via t-distribution ($t_{0.025,2}=4.303$). Promoter results: seed 42 = 0.890, seed 1337 = 0.878, seed 2026 = 0.891, mean $0.886\pm0.007$ (95\% CI $\pm0.018$). All TF-binding tasks show std $\leq 0.019$, confirming stable performance ordering. Splice-site F1 is identical across all seeds (0.241), indicating the failure mode is deterministic.}
    \label{tab:multiseed}
    \begin{tabular}{lrrrrc}
        \toprule
        Task & Seed 42 & Seed 1337 & Seed 2026 & Mean & 95\% CI \\
        \midrule
        human\_tf\_0  & 0.791 & 0.807 & 0.770 & 0.789 & $\pm 0.047$ \\
        human\_tf\_1  & 0.804 & 0.807 & 0.823 & 0.811 & $\pm 0.026$ \\
        human\_tf\_2  & 0.718 & 0.710 & 0.724 & 0.717 & $\pm 0.017$ \\
        human\_tf\_3  & 0.619 & 0.605 & 0.546 & 0.590 & $\pm 0.096$ \\
        human\_tf\_4  & 0.810 & 0.808 & 0.805 & 0.808 & $\pm 0.007$ \\
        splice\_site  & 0.241 & 0.241 & 0.241 & 0.241 & $\pm 0.000$ \\
        \bottomrule
    \end{tabular}
\end{table}

\begin{table}[H]
    \centering
    \small
    \caption{\textbf{Training Budget Parity.} 20-epoch fine-tuned \dnabert baselines across all GUE tasks. Only human\_tf\_1 showed strong underfitting at 5 epochs (0.500 → 0.872 at 20 epochs). All other tasks change by $<$0.02 F1, confirming that the 5-epoch budget is adequate for those tasks and that the human\_tf\_1 anomaly is an isolated case. See also Table~\ref{tab:multiseed} for RL multi-seed results.}
    \label{tab:budget_parity}
    \begin{tabular}{lrrrr}
        \toprule
        Task & 5-epoch F1 & 20-epoch F1 & $\Delta$F1 & Verdict \\
        \midrule
        human\_tf\_1  & 0.333 & 0.872 & +0.539 & Underfitting \\
        human\_tf\_0  & 0.842 & 0.839 & $-$0.003 & Adequate \\
        human\_tf\_2  & 0.838 & 0.833 & $-$0.005 & Adequate \\
        human\_tf\_3  & 0.790 & 0.791 & +0.001 & Adequate \\
        human\_tf\_4  & 0.887 & 0.894 & +0.007 & Adequate \\
        splice\_site  & 0.924 & 0.923 & $-$0.001 & Adequate \\
        \bottomrule
    \end{tabular}
\end{table}

\begin{table}[H]
    \centering
    \small
    \caption{Decoding threshold robustness for fine-tuned RL grouping on promoter detection. The downstream model is fully re-run from scratch for each threshold with the new segmentation. Accuracy is stable across thresholds because mean-pooling preserves global sequence information even with few segments: at $t{=}0.50$ (1 token), the single segment's representation equals the mean of all 63 token states, which retains global promoter statistics sufficient for binary classification. The threshold affects which local patterns are captured, not whether global information is available.}
    \label{tab:threshold}
    \begin{tabular}{rrrr}
        \toprule
        Threshold & Acc & F1 & Comp \\
        \midrule
        0.10 & 0.893 & 0.893 & 0.973 \\
        0.12 & 0.891 & 0.890 & 0.732 \\
        0.15 & 0.889 & 0.889 & 0.483 \\
        0.20 & 0.891 & 0.891 & 0.193 \\
        0.25 & 0.892 & 0.892 & 0.132 \\
        0.30 & 0.893 & 0.893 & 0.089 \\
        0.40 & 0.893 & 0.893 & 0.027 \\
        \bottomrule
    \end{tabular}
\end{table}

\begin{table}[H]
    \centering
    \caption{Entropy-bonus ablation on frozen-env promoter detection.}
    \label{tab:entropy}
    \begin{tabular}{lrrr}
        \toprule
        Method & Acc & Tok & $\comp$ \\
        \midrule
        Frozen RL, $\beta=0.02$ (baseline)   & 0.7546 & 4.22 & 0.0665 \\
        Frozen RL + entropy ($w=0.05$)        & 0.7549 & 5.04 & 0.0839 \\
        \bottomrule
    \end{tabular}
\end{table}%
\section{Additional Figures}%
\begin{figure*}[!h]
    \centering
    \includegraphics[width=\linewidth]{figures/generated/fig_training_dynamics.pdf}
    \caption{Training dynamics across three seeds. Classifier loss decreases steadily while reward and token count stabilize, confirming consistent learning without collapse.}
    \label{fig:supp_training}
\end{figure*}

\begin{figure*}[!h]
    \centering
    \includegraphics[width=\linewidth]{figures/generated/fig_token_length_distribution.pdf}
    \caption{Token length distributions for single-nucleotide, fixed 6-mer, and learned RL tokenizers. Learned RL achieves similar average compression to fixed 6-mer ($\sim$0.18) but with a variable-length distribution reflecting context-dependent segmentation.}
    \label{fig:supp_toklen}
\end{figure*}

\begin{figure*}[!ht]
    \centering
    \includegraphics[width=\linewidth]{figures/generated/fig_expanded_motif_rates.pdf}
    \caption{Expanded motif boundary enrichment across six regulatory motifs. TATA shows strong on-motif enrichment ($\sim$5$\times$); CAAT is marginal; GCGC, SP1, E-box, and GC-box all fall at or below the off-motif baseline. TATA enrichment is a specific learned pattern rather than a general response to GC content.}
    \label{fig:supp_expanded_motifs}
\end{figure*}

\begin{figure*}[!ht]
    \centering
    \includegraphics[width=\linewidth]{figures/generated/fig_positionwise_heatmap.pdf}
    \caption{Position-wise mean boundary probability (top) and hard boundary rate (bottom) across DNABERT-2 token positions on the test set. Positive and negative promoters show broadly similar profiles with no strong positional bias, suggesting learned boundaries are distributed throughout the sequence.}
    \label{fig:supp_positionwise}
\end{figure*}
\section{Appendix: Threshold Invariance and Compression Control}
\begin{table}[!ht]
\centering
\small
\begin{tabular}{ccccc}
\toprule
\textbf{Threshold ($\tau$)} & \textbf{Accuracy} & \textbf{F1} & \textbf{Compression} & \textbf{Avg Tokens} \\
\midrule
0.05 & 0.8929 & 0.8929 & 1.0000 & 63.4 \\
0.08 & 0.8929 & 0.8929 & 1.0000 & 63.4 \\
0.10 & 0.8932 & 0.8932 & 0.9726 & 61.7 \\
0.12 & 0.8905 & 0.8904 & 0.7320 & 46.4 \\
0.15 & 0.8894 & 0.8893 & 0.4830 & 30.6 \\
0.20 & 0.8910 & 0.8910 & 0.1925 & 12.2 \\
0.25 & 0.8921 & 0.8920 & 0.1315 & 8.34 \\
0.30 & 0.8927 & 0.8927 & 0.0892 & 5.66 \\
0.40 & 0.8929 & 0.8929 & 0.0271 & 1.72 \\
0.50 & 0.8929 & 0.8929 & 0.0158 & 1.00 \\
\bottomrule
\end{tabular}
\caption{\textbf{Threshold Robustness Analysis (Promoter Task).} RL boundary policy threshold invariance demonstrates that accuracy remains stable (±0.003) across a 63.4× variation in compression (0.016 → 1.0). Most notably, at $\tau=0.50$ (single segment, 1 token), accuracy = 0.8929 and macro-F1 = 0.8929 — identical values. Since all GUE test splits are perfectly balanced (50\% positive, 50\% negative), a majority-class predictor would achieve accuracy = 0.500 and F1 = 0.333 (predicting only one class). The convergence of accuracy and macro-F1 at 0.8929 therefore rules out degenerate majority-class prediction and confirms that the model genuinely classifies from the global mean representation. The single-segment embedding (arithmetic mean of all 63 DNABERT-2 token states) preserves sufficient global sequence statistics for binary classification: GC content, dinucleotide composition, and mean contextual embedding direction are all captured in the global average. The classifier's decision boundary exploits these global statistics rather than local boundary cues, which explains why learned boundary placement has minimal impact on aggregate accuracy. This finding supports the conclusion that compression-based regularization drives performance, while learned boundaries add value primarily through biological interpretability (TATA enrichment) rather than aggregate classification.}
\label{tab:threshold_robustness}
\end{table}%
\begin{table}[!ht]
\centering
\small
\begin{tabular}{lcccc}
\toprule
\textbf{Task} & \textbf{Target Compression} & \textbf{Achieved Compression} & \textbf{Variance (\%)} & \textbf{F1} \\
\midrule
promoter & 0.1800 & 0.1222 & -32.1 & 0.8123 \\
human\_tf\_0 & 0.1800 & 0.1417 & -21.3 & 0.8359 \\
human\_tf\_1 & 0.1800 & 0.1389 & -22.8 & 0.8677 \\
human\_tf\_2 & 0.1800 & 0.1371 & -23.9 & 0.7434 \\
human\_tf\_3 & 0.1800 & 0.1422 & -21.0 & 0.6083 \\
human\_tf\_4 & 0.1800 & 0.1424 & -20.9 & 0.7921 \\
splice\_site & 0.1800 & 0.1254 & -30.3 & 0.2453 \\
\bottomrule
\end{tabular}
\caption{\textbf{Lagrangian RL Compression Control.} Hard dual-ascent constraints successfully enforce tight compression budgets across all 7 GUE tasks. Achieved compression averages 0.1357 (23.7\% below target of 0.18), with per-task variance in the range $-21\%$ to $-32\%$. All tasks cluster tightly around the target, demonstrating that hard constraints are more effective than soft $\beta$-penalties for compression control. While Lagrangian RL improves F1 on TF-binding tasks (avg +2.7-6.4\%), splice-site shows minimal improvement (+1.8\%), indicating that architectural limitations (loss of positional signal via mean-pooling) persist despite improved constraint enforcement.}
\label{tab:lagrangian_compression}
\end{table}%
\begin{table}[!ht]
\centering
\small
\begin{tabular}{lccc}
\toprule
\textbf{Task \& Method} & \textbf{Compression} & \textbf{F1} & \textbf{Accuracy} \\
\midrule
\multicolumn{4}{l}{\textit{Promoter Detection}} \\
Fine-tuned DNABERT-2 & 1.000 & 0.939 & 0.939 \\
RL (fine-tuned) & 0.190 & 0.886 & 0.890 \\
In-encoder token merging & 0.244 & 0.869 & 0.869 \\
Stride pooling ($s=4$) & 0.210 & 0.920 & 0.920 \\
Gumbel-Softmax & 0.210 & 0.876 & 0.876 \\
\midrule
\multicolumn{4}{l}{\textit{Splice-Site Prediction}} \\
DNABERT-2 & 1.000 & 0.929 & 0.929 \\
RL (fine-tuned) & 0.250 & 0.241 & 0.566 \\
In-encoder token merging & 0.245 & 0.146 & 0.225 \\
Stride pooling ($s=4$) & 0.255 & 0.925 & 0.925 \\
\bottomrule
\end{tabular}
\caption{\textbf{Baseline Comparison at Matched Compression Ratios.} RL-learned boundaries versus fixed-stride pooling, Gumbel-Softmax segmentation, and in-encoder token merging (ToMe-style bipartite similarity merging on DNABERT-2 final hidden states) at comparable compression budgets. On promoter detection, fine-tuned RL (F1: 0.886) outperforms in-encoder token merging (F1: 0.869) at matched compression, demonstrating that learned pre-encoding boundaries provide a measurable advantage over post-hoc similarity-based pruning. However, fixed-stride pooling (F1: 0.920) outperforms both, indicating that uniform compression remains the strongest baseline. On splice-site, all compression methods fail relative to the uncompressed baseline, but RL (F1: 0.241) substantially outperforms token merging (F1: 0.146). RL's distinctive value lies in biological interpretability (TATA enrichment, Sec.~\ref{sec:interpretability}) rather than aggregate predictive accuracy.}
\label{tab:baseline_compression}
\end{table}

\begin{table}[!ht]
\centering
\small
\begin{tabular}{llrrrr}
\toprule
\textbf{Model} & \textbf{Class} & \textbf{Precision} & \textbf{Recall} & \textbf{F1} & \textbf{ECE} \\
\midrule
\multirow{3}{*}{FT \dnabert baseline} & no-splice & 0.916 & 0.924 & 0.920 & \multirow{3}{*}{0.048} \\
 & acceptor & 0.921 & 0.907 & 0.914 & \\
 & donor & 0.936 & 0.939 & 0.938 & \\
\midrule
\multirow{3}{*}{Fine-tuned RL} & no-splice & 0.000 & 0.000 & 0.000 & \multirow{3}{*}{0.024} \\
 & acceptor & 0.000 & 0.000 & 0.000 & \\
 & donor & 0.566 & 1.000 & 0.723 & \\
\bottomrule
\end{tabular}
\caption{\textbf{Per-class splice-site analysis.} Fine-tuned RL collapses to predicting only the majority ``donor'' class (recall = 1.0, all other classes recall = 0.0), fully explaining the accuracy/F1 divergence (0.566/0.241). The confusion matrix confirms all 956 no-splice and 1025 acceptor test examples are misclassified as donor. Low ECE (0.024) reflects confident but completely biased predictions. The fine-tuned \dnabert baseline achieves balanced per-class F1 across all three classes. Five reward-shaping variants (curriculum, asymmetric recall, GT/AG proximity $\gamma \in \{0.1, 0.3\}$) all reproduce the same collapse, confirming the failure is representational rather than reward-related.}
\label{tab:splice_perclass}
\end{table}

\end{document}
